The arguments of Sections~\ref{sec:quotients} and~\ref{sec:higgs} use the spacetime dimension only
through the supersymmetry conditions~\eqref{eq:vacua}. In this section we impose the corresponding
conditions of four-dimensional $N=2$ supergravity directly. The resulting statement concerns type~IIA
string theory on $\widehat X$ and, through its hypermultiplet moduli space, the setting of Greene,
Morrison and Vafa~\cite[Sect.~6]{Greene:1996Confinement}. Type~IIA on $\widehat X$ is M-theory on
$\widehat X\times S^1$, but the argument does not pass through the five-dimensional theory: it is carried
out directly in four dimensions, where the massless states at the root carry no D0-brane charge, so that
no tower of Kaluza--Klein (D0) states has to be truncated.

\subsection{Supersymmetric Minkowski vacua in four dimensions}
\label{subsec:4d-setting}

We use the conventions of~\cite{Andrianopoli:1996cm,Hristov:2009}: $n_V$ vector multiplets with special
K\"ahler manifold $\mathcal S$, homogeneous coordinates $X^\Lambda$, $\Lambda=0,\dots,n_V$, a prepotential
$\mathcal F(X)$, symplectic sections $(\mathcal L^\Lambda,\mathcal M_\Lambda)=e^{K/2}(X^\Lambda,\partial_\Lambda\mathcal F)$,
$f_i^\Lambda=\nabla_i\mathcal L^\Lambda$ and metric $g_{i\bar\jmath}$; hypermultiplets with quaternionic
K\"ahler manifold $M$ of reduced scalar curvature $\nu<0$ fixed by supergravity; and an electric gauging of
a torus acting on $M$ by isometries, with Killing vectors $k_\Lambda$ and triplets of moment maps
$\vec P_\Lambda$, and acting trivially on $\mathcal S$. For a light state of charge $q\in\ZZ^{n_V+1}$ we write
$Z_q=e^{K/2}q_\Lambda X^\Lambda$ for its central charge.

For configurations with constant scalars and vanishing field strengths the gravitino, gaugino and
hyperino variations vanish if and only if $\vec P_\Lambda\mathcal L^\Lambda=0$,
$g^{i\bar\jmath}\bar f^\Lambda_{\bar\jmath}\vec P_\Lambda=0$ and $k_\Lambda\mathcal L^\Lambda=0$~\cite[Sects.~2.1--2.3,
2.4.2]{Hristov:2009}. Since the $\vec P_\Lambda$ are real and, where $g_{i\bar\jmath}$ is positive definite,
the $(n_V+1)\times(n_V+1)$ matrix $(\bar{\mathcal L}^\Lambda,f_i^\Lambda)$ is
invertible~\cite[Lemma~B.5, Cor.~B.6]{crapsetal1997specialgeometry}, the conditions become
\begin{equation}
 \vec P_\Lambda=0\quad\text{for all }\Lambda,\qquad k_\Lambda\mathcal L^\Lambda=0,
 \label{eq:vacua-4d}
\end{equation}
as stated in~\cite[Sect.~2.4.2]{Hristov:2009}. Positive definiteness of $g_{i\bar\jmath}$ is equivalent to
negative definiteness of the imaginary part of the gauge kinetic matrix
$\mathcal N_{\Lambda\Sigma}$~\cite[eq.~(4.45)]{crapsetal1997specialgeometry}; it takes the place of
condition~(S3). We write $n_N$ for the quaternionic dimension of the manifold $N$ of neutral hypermultiplets. Given charge
vectors $q_1,\dots,q_k$ of $k$ charged hypermultiplets, a \emph{root point} is a point $x=(z_0,b,0)$ with
$z_0\in\mathcal S$ at which $Z_{q_1},\dots,Z_{q_k}$ vanish, $b\in N$, and vanishing charged hypermultiplets;
$\mathcal V_x$ is defined as in Definition~\ref{def:vacuum-germ}, with $z_0$ in place of $\phi$ and the theory
of the following hypothesis.

\begin{hypothesis}[regularity, four dimensions]\label{hyp:regular-4d}
Near a root point $x=(z_0,b,0)$ the dynamics below the lightest massive state is a four-dimensional $N=2$
gauged supergravity as described above, in which exactly $k$ charged hypermultiplets are light, with: $g_{i\bar\jmath}$ positive definite at
$z_0$; a smooth quaternionic K\"ahler manifold $M$ of quaternionic dimension $n_N+k$ containing the
manifold $N$ of neutral hypermultiplets, including the corrections to its metric from D2-instantons on
three-cycles that do not vanish at the root and from NS5-instantons; charge vectors
$q_1,\dots,q_k$ of the light states spanning a lattice of rank $k-1$ with the single relation
$\sum_ia_iq_i=0$, $a$ primitive and every $a_i\neq0$; no neutral hypermultiplet charged; and $b$ a smooth
point of $N$. Moreover, the Coulomb branch through $x$ consists of supersymmetric Minkowski vacua, and the
higher-derivative couplings do not change the set of such vacua.
\end{hypothesis}

As in five dimensions (Section~\ref{subsec:other-effects}), D2-instantons on three-cycles that do not
vanish and NS5-instantons correct the metric of $M$ smoothly and preserve the gauged isometries. We expect, as in
the three-dimensional analysis of~\cite{SeibergShenker:1996}, that the Euclidean D2-branes on the vanishing
three-spheres, which produce the singularity of the neutral metric in~\cite{OoguriVafa:1996}, describe the
same singular point in terms of the neutral fields alone: in the Wilsonian description their effect is
carried by the light charged hypermultiplets and, through the massive D2--D0 states, by smooth corrections
to $M$. This is part of the content of the identification (W-IIA) of Section~\ref{subsec:gmv}.

\begin{theorem}\label{thm:4d}
Assume Hypothesis~\ref{hyp:regular-4d}, write $Z_i=Z_{q_i}$, and let $\mathcal S_0=\{Z_1=\dots=Z_k=0\}$.
Then $\mathcal S_0$ is a complex submanifold of $\mathcal S$ of codimension $k-1$ near $z_0$ and, writing
$(\mathcal S,z_0)$ and $(\mathcal S_0,z_0)$ for germs, as germs of stratified spaces
\begin{equation}
 \mathcal V_x\;\cong\;\bigl((\mathcal S,z_0)\times N\bigr)\;\cup\;\bigl((\mathcal S_0,z_0)\times N\times\CC^2/\ZZ_m\bigr),
 \qquad m=\sum_i|a_i|,
 \label{eq:thm-4d}
\end{equation}
glued along $(\mathcal S_0,z_0)\times N\times\{0\}$. There is no mixed branch; the hypermultiplet
factor of the Higgs branch off its root is a smooth quaternionic K\"ahler manifold of quaternionic
dimension $n_N+1$, and at distance $r$ from the root, measured in the metric of $M$ normalised as
in~\cite{Hristov:2009} ($\kappa^2=1$), its metric is the
flat orbifold metric of $\CC^2/\ZZ_m$ times the flat metric on $T_bN$ up to relative corrections of order
$r^2$.
\end{theorem}

\begin{proof}
In a frame with a prepotential the matrix $(X^\Lambda,\partial_jX^\Lambda)$ is
invertible~\cite[Sect.~4, on the existence of a prepotential]{crapsetal1997specialgeometry}, so a
combination $\sum_ic^iq_i$ annihilating both $X$ and $\partial X$ at a point of $\{q\cdot X=0\}$ vanishes;
since $\sum_ia_iZ_i\equiv0$ and the $q_i$ have rank $k-1$, $\mathcal S_0$ has codimension $k-1$. On the
Coulomb branch through $x$ the conditions~\eqref{eq:vacua-4d} hold with the charged hypermultiplets set
to zero, so the moment map vanishes at the point $p\in M$ over $b$, and on $N$ near $b$ by
Lemma~\ref{lem:fixed}. The torus acting effectively on the charged hypermultiplets is the image of
$U(1)^{n_V+1}$ in $U(1)^k$, whose cocharacter lattice is $a^\perp\cap\ZZ^k$, so Lemma~\ref{lem:flat-model}
applies without a saturation assumption. Section~\ref{sec:quotients} and Lemma~\ref{lem:flat-model} do
not refer to the spacetime dimension, and Theorem~\ref{thm:normal-form},
Proposition~\ref{prop:tangent-cone} and Lemma~\ref{lem:flat-model} give the hypermultiplet germ
$N\times\CC^2/\ZZ_m$ and the metric statement as in the proof of Theorem~\ref{thm:general}. For the
vector conditions, $k_\Lambda\mathcal L^\Lambda$ is the fundamental field of the element $\xi$ of
$\mathrm{Lie}(T)\otimes\CC$ determined by $(Z_1,\dots,Z_k)$. Off the root the effective torus acts freely
near $p$, so the real and imaginary parts of $\xi$ vanish, that is $Z_i=0$ for all $i$; conversely $Z=0$
gives $k_\Lambda\mathcal L^\Lambda=0$. On the Coulomb branch every $z$ is allowed. The absence of a mixed
branch follows as in Theorem~\ref{thm:general}.
\end{proof}

The Higgs branch has vector factor $(\mathcal S_0,z_0)$ of complex dimension $n_V-(k-1)$. For type~IIA on
$\widehat X$ at large volume (Section~\ref{subsec:iia}), with the D0-charges fixed by the gauge choice
$B\cdot C_i=0$ made there, the root points with $J\in F$ are the points of $\mathcal E_F\times N$,
and $n_V-(k-1)=h^{1,1}(X_t)$, the number of vector multiplets of type~IIA on the smoothing.

\subsection{Type IIA string theory at the massless locus}
\label{subsec:iia}

Let $\widehat X$ and $F$ satisfy (S1)--(S3). We use the large-volume symplectic frame of type~IIA on
$\widehat X$, with complexified K\"ahler moduli $w=B+iJ$, where $w^A=X^A/X^0$ and the prepotential is
$\mathcal F=-\tfrac16\kappa_{ABC}X^AX^BX^C/X^0$ plus worldsheet-instanton terms. In this frame a D2-brane on a class
$\beta$ bound to $n$ D0-branes has holomorphic central charge $\beta\cdot w-n$, up to the factor $e^{K/2}$
and the sign conventions of~\cite[eq.~(A.10)]{Denef:2007Halos}; D2- and D0-charges are
electric~\cite{Denef:2007Halos}. The vector multiplets gauge no isometry of $\mathcal S$. We call
\[
 \mathcal E_F=\{\,w=B+iJ:\ J\in F,\ B\cdot C_i\in\ZZ\ \text{for all }i\,\}
\]
the \emph{massless locus} over $F$, the locus of the complexified K\"ahler cone at which D2-branes on the
curves $C_i$ become massless. The large gauge transformations $B\mapsto B+D$, $D\in H^2(\widehat X;\ZZ)$,
identify its components and shift the D0-charges; since the charge lattice is saturated by (S2), we can
use them to set $B\cdot C_i=0$. A D2-brane on a class
in the span of the $C_i$ then has central charge proportional to $n$, so among these states only the
D2-branes on the $C_i$ with $n=0$ are massless; by (S1) and the vanishing of multiple-cover and
higher-genus invariants of isolated $(-1,-1)$-curves~\cite{FaberPandharipande:2000,BryanKatzLeung:2001}
they give $k$ hypermultiplets with charges $q_{iA}=H_A\cdot C_i$, $q_{i0}=0$. D4- and D6-brane central
charges receive worldsheet-instanton corrections and can vanish at finite volume, for instance for
D4-branes on a small divisor near the del Pezzo facet of $X_9$. We therefore work at large volume,
$J=\rho J_0$ with $J_0$ in a compact subset of $F$ and $\rho\geq\rho_0(J_0)$, and assume, as part of
Hypothesis~\ref{hyp:regular-4d}, that there the central charges of all other BPS states are bounded away
from zero (for D0-branes and D2--D0 states on the span of the $C_i$ this holds exactly, $|n|\geq1$);
$\rho_0(J_0)\to\infty$ as $J_0$ approaches the boundary of $F$.

\begin{proposition}\label{prop:4d-kinetic}
For the cubic prepotential, with $L(J)_{AB}=\kappa_{ABC}J^C$, $v=L(J)J$ and $\operatorname{vol}=\tfrac16\kappa(J,J,J)$,
one has $e^{-K}=\tfrac43\kappa(J,J,J)$ and the special K\"ahler metric on the complexified K\"ahler cone is
\begin{equation}
 g_{A\bar B}=\frac{1}{4\operatorname{vol}}\,\mathcal G_{AB},\qquad
 \mathcal G=-L(J)+\frac{vv^{\mathsf T}}{4\operatorname{vol}},
 \label{eq:4d-kinetic}
\end{equation}
where $\mathcal G$ is the Hodge metric on harmonic two-forms of (S3), that is the five-dimensional gauge
kinetic matrix~\cite[eqs.~(4.1), (4.6)]{W26quantum}. In particular, on $F$ the metric of the cubic
prepotential is positive definite if and only if (S3) holds.
\end{proposition}

\begin{proof}
With $X^0=1$, $e^{-K}=i(\bar X^\Lambda\partial_\Lambda\mathcal F-X^\Lambda\overline{\partial_\Lambda\mathcal F})
=\tfrac43\kappa(J,J,J)$. Since $K$ depends on $J$ only,
$\partial_{w^A}\partial_{\bar w^B}K=\tfrac14\partial_{J^A}\partial_{J^B}K$, and with
$\partial_A\kappa(J,J,J)=3v_A$, $\partial_A\partial_B\kappa(J,J,J)=6L_{AB}$ one obtains
$g=\tfrac14\bigl(-6L/\kappa(J,J,J)+9vv^{\mathsf T}/\kappa(J,J,J)^2\bigr)$, which
is~\eqref{eq:4d-kinetic} with $\kappa(J,J,J)=6\operatorname{vol}$.
\end{proof}

The corrections to the Wilsonian prepotential near $\mathcal E_F$ are the constant $\zeta(3)\chi$ term, of
relative order $\rho^{-3}$; worldsheet instantons on classes of positive area, exponentially small; and the
regular part of $\operatorname{Li}_3(e^{2\pi i\,w\cdot C_i})$ at the shrinking curves, whose non-analytic part,
proportional to $(w\cdot C_i)^2\log(w\cdot C_i)$, is the one-loop contribution of the light hypermultiplets
and is absent from the Wilsonian prepotential. The regular part changes $g$ by a relative amount of order
$\rho^{-1}$ in the directions normal to $\mathcal S_0$, with a coefficient fixed by the Wilsonian
subtraction scale; for any fixed subtraction scale, the positive definiteness in
Hypothesis~\ref{hyp:regular-4d} follows from (S3) for $\rho$ large. Proposition~\ref{prop:face}(3)
establishes (S3) for $X_9$.

\subsection{The vanishing three-spheres}
\label{subsec:spheres}

\begin{lemma}\label{lem:spheres}
Let $X_t$ be a smoothing of the contraction $X_F$ of Section~\ref{subsec:setting}, and $S_i\subset
X_t$ the vanishing three-sphere of the $i$-th node. For suitable orientations of the $S_i$,
$[S_i]=a_i[S]$ in $H_3(X_t;\ZZ)$ for a class $[S]$ generating the span of the $[S_i]$. If the relation
$\sum_ia_i[C_i]=0$ holds in $H_2(\widehat X;\ZZ)$ and not only modulo torsion, then $[S]$ is primitive.
\end{lemma}

\begin{proof}
The complements $\widehat X\smallsetminus\bigcup_iC_i$ and $X_t\smallsetminus\bigcup_iS_i$ are
diffeomorphic, and the linking three-sphere of $C_i$ is identified with $S_i$. In the exact sequence
$H_4(\widehat X)\to\bigoplus_iH_4(\widehat X,\widehat X\smallsetminus C_i)\cong\ZZ^k\to
H_3(\widehat X\smallsetminus\textstyle\bigcup C_i)$ the first map is $D\mapsto(D\cdot C_i)_i$, whose
image is $a^\perp\cap\ZZ^k$ by (S2), and the second sends the $i$-th generator to the linking
sphere. The map $H_3(X_t\smallsetminus\bigcup S_i)\to H_3(X_t)$ is injective, because
$H_4(X_t,X_t\smallsetminus\bigcup S_i)\cong H^2(\bigsqcup S_i)=0$ by Lefschetz duality. So the span of
the $[S_i]$ is $\ZZ^k/(a^\perp\cap\ZZ^k)\cong\ZZ$, via $c\mapsto a\cdot c$. For primitivity consider
$H_3(X_t)\to H_3(X_t,X_t\smallsetminus\bigcup S_i)\cong\ZZ^k\to H_2(X_t\smallsetminus\bigcup S_i)\cong
H_2(\widehat X\smallsetminus\bigcup C_i)\cong H_2(\widehat X)$, the last isomorphism because
$H_j(\widehat X,\widehat X\smallsetminus\bigcup C_i)\cong H^{6-j}(\bigsqcup C_i)=0$ for $j=2,3$. The first map
is $\gamma\mapsto(\gamma\cdot S_i)_i$ and the second sends the $i$-th generator to $\pm[C_i]$, so the image
of the first is the integral relation lattice of the curves, $\ZZ a$. Hence some $\gamma$ has
$\gamma\cdot S=\pm1$, and $[S]$ is primitive.
\end{proof}

For unit coefficients the $k$ vanishing spheres are homologous up to orientation: this is the
configuration of $k$ homologous vanishing three-cycles of~\cite{OoguriVafa:1996,Greene:1996Confinement}.
For $X_9$, Lemma~\ref{lem:spheres} with~\eqref{eq:face-relation} gives four such spheres.

\subsection{The expectation of Greene, Morrison and Vafa}
\label{subsec:gmv}

Greene, Morrison and Vafa argued from the rigid field theory that near homologous vanishing three-cycles
the type~IIA hypermultiplet moduli space has, with gravity decoupled, an $A_{k-1}$ singularity, $k$ the
number of cycles, and wrote: ``When we turn on gravity we do expect that the qualitative features we have
found will persist but the actual metric will get corrected in a non-universal way, depending on how the
vanishing cycle sits in the Calabi--Yau, and the total space will be a quaternionic
manifold''~\cite[Sect.~6]{Greene:1996Confinement}. Here ``gravity decoupled'' is the limit in which the
four-dimensional dilaton, a coordinate of the universal hypermultiplet, is sent to weak coupling while
approaching the conifold locus, so that the relevant piece of the quaternionic K\"ahler metric scales to a
hyperk\"ahler one~\cite[Sect.~6]{Greene:1996Confinement}; ``finite Planck mass'' means the quaternionic
K\"ahler metric at fixed $\nu$.

Under (W-IIA) below, we identify the neutral manifold $N$ of $\widehat X$ with the locus in the hypermultiplet
moduli space of the smoothings $X_t$ where they degenerate to $X_F$: its points are the complex structures of $X_F$ with their
Ramond--Ramond axions and the scalars of the universal hypermultiplet. We call $N\times\{0\}$ the \emph{root locus}; it lies in the conifold
locus of $\mathcal M_H(X_t)$, where the vanishing cycles have zero period. To compare the vacuum space of
Theorem~\ref{thm:4d} with the moduli space of type~IIA on $X_t$ we assume, at large volume:
\begin{itemize}
\item[(W-IIA)] near $x$, the Higgs vacua off the root are type~IIA compactifications on the smoothings
$X_t$, and the hypermultiplet moduli space $\mathcal M_H^{\mathrm{reg}}$ of these (the same manifold as in
Section~\ref{subsec:object}, by~\cite[Sect.~2]{Strominger:1997}) is bi-Lipschitz
equivalent near the root to $(\mu^{-1}(0)\smallsetminus N)/T$ with the quotient metric
of~\cite{Galicki:1988}.
\end{itemize}
This is the four-dimensional counterpart of clause~(3) of Hypothesis~\ref{hyp:regular}, and it is the
assumption on which the comparison with~\cite{Greene:1996Confinement} rests. Without gravity its analogue
holds for the known model metrics: the hyperk\"ahler metrics obtained from Euclidean D2-branes on $k$
homologous vanishing spheres~\cite{OoguriVafa:1996,SeibergShenker:1996} (constructed rigorously for $k=1$
in~\cite{FredricksonZimet:2025}, whose Remarks~3.14 and~4.44 indicate the extension to multiplicity $k$)
have the same $A_{k-1}$ point as the hyperk\"ahler quotient of the $k$ light
hypermultiplets. A proof at finite Planck mass would require control of the instanton-corrected
quaternionic K\"ahler metric near the conifold locus. The twistorial construction~\cite{Alexandrov:2008gh} describes the D-instanton corrections to
this metric, NS5-instantons aside, through the invariants of the wrapped branes, but its behaviour where the central charge of a wrapped brane vanishes has not been
analysed there, and it does not settle (W-IIA).

\begin{corollary}\label{cor:gmv}
Let $\widehat X$ and $F$ satisfy (S1)--(S3), and assume Hypothesis~\ref{hyp:regular-4d} and (W-IIA) for
type~IIA on $\widehat X$ at large volume. At a point of the root locus the germ of the metric completion of
the hypermultiplet moduli space of $X_t$ is homeomorphic to $N\times\CC^2/\ZZ_m$ with $m=\sum_i|a_i|$, and
it is a smooth quaternionic K\"ahler manifold off $N\times\{0\}$. If the equivalence in (W-IIA) is an
isometry, the metric is the flat orbifold metric up to relative corrections of order $r^2$.
\end{corollary}

\begin{proof}
The argument of Lemma~\ref{lem:completion} for the hypermultiplet factor, with
Theorem~\ref{thm:4d} in place of Theorem~\ref{thm:general} and (W-IIA) in place of clause (3) of
Hypothesis~\ref{hyp:regular}; the metric clause follows from Proposition~\ref{prop:tangent-cone} when the
equivalence is an isometry.
\end{proof}

For unit coefficients $m=k$, the number of homologous vanishing spheres. Near the root the rest of the
conifold locus consists of smooth points of the completion: in $\CC[X,Y,s]/(XY-c\,s^m)$, $c\neq0$ as in Lemma~\ref{lem:flat-model}, it is
$\{s=0,\,(X,Y)\neq0\}$. For coefficients $|a_i|\geq2$ the rigid instanton picture is consistent with this description of the root.
Near the conifold locus the rigid metric has Gibbons--Hawking form, with coordinates the period $\int_S\Omega$
and the Ramond--Ramond phase $\theta=\int_SC_3$ and a harmonic potential. A BPS state of charge $n$ with
respect to $[S]$ (for the sphere $S_i$, $n=|a_i|$) in the three-dimensional description obtained by compactification on a
circle~\cite{SeibergShenker:1996} contributes $n$ singular points to this
potential, at $\theta\in2\pi\ZZ/n$, each of which is an $A_{n-1}$ point~\cite[Sect.~4.1]{Gaiotto:2008cd}, and
the one-loop potential is a sum over the light hypermultiplets~\cite{SeibergShenker:1996,Gaiotto:2008cd}.
Hence at $\theta=0$, where the centres of all the $S_i$ coincide, the singularity is of type $A_{m-1}$ with
$m=\sum_i|a_i|$. The rigid picture also places orbifold points
at other values of $\theta$, outside the germ; we make no statement about them.
